% Part IV -- Numerical Data, Tables, and Visualization
% Chapter 11: NumPy arrays and vectorized calculations

\h{NumPy Arrays and Vectorized Calculations}

Python lists work well when a collection contains mixed kinds of values. When
the values are mainly numbers and share one structure, a NumPy array makes
calculations shorter and easier to check. This chapter moves from one row of
numbers to small two-dimensional tables.

\begin{NoteBox}
\textbf{Prerequisites.} You should be comfortable with variables, lists,
indexing, loops, functions, comparisons, and copying mutable data.
\end{NoteBox}

\begin{tcolorbox}[title=Learning outcomes,colframe=orange,colback=white,arc=2mm]
After completing this chapter, you should be able to:
\begin{itemize}
    \item create one- and two-dimensional NumPy arrays;
    \item inspect an array's shape, size, number of dimensions, and data type;
    \item select values with indices, slices, and boolean masks;
    \item calculate with whole arrays instead of writing a loop for each value;
    \item summarize all values or summarize along one axis;
    \item reshape compatible data and use simple broadcasting; and
    \item make an independent copy before changing selected data.
\end{itemize}
\end{tcolorbox}

\hh{Install once and import when needed}

NumPy is an external library. Install it once in the environment where you run
Python. The command below belongs in a terminal, not inside a Python file.

\begin{verbatim}
python -m pip install numpy
\end{verbatim}

Import NumPy in every program that uses it. The short name \c{np} is the
standard convention.

\code{c11_import_numpy}{python}{
    import numpy as np

    values = np.array([12.0, 15.5, 18.0])
    print(values)
}{Import NumPy and create a first array}

If the import fails, check that NumPy was installed for the same Python
environment that runs the file.

\hh{Create an array from a list}

The \c{np.array()} function converts a sequence into an array. A one-dimensional
array behaves like one ordered row of values.

\code{c11_array_from_list}{python}{
    import numpy as np

    room_areas_m2 = np.array([24.0, 18.5, 31.0, 22.5])

    print(room_areas_m2)
    print(type(room_areas_m2))
}{Create a one-dimensional array}

An array normally uses one data type for all its elements. This regular
structure is one reason NumPy can process numerical data efficiently.

\hh{Inspect an array before using it}

Four attributes answer different questions about an array.

\begin{center}
\begin{tabular}{ll}
\textbf{Attribute} & \textbf{Question answered} \\
\c{ndim} & How many axes does the array have? \\
\c{shape} & How many positions are on each axis? \\
\c{size} & How many values are stored in total? \\
\c{dtype} & What data type stores the values? \\
\end{tabular}
\end{center}

\code{c11_array_attributes}{python}{
    import numpy as np

    room_areas_m2 = np.array([24.0, 18.5, 31.0, 22.5])

    print("Dimensions:", room_areas_m2.ndim)
    print("Shape:", room_areas_m2.shape)
    print("Size:", room_areas_m2.size)
    print("Data type:", room_areas_m2.dtype)
}{Inspect a one-dimensional array}

For this array, \c{shape} is \c{(4,)}. The comma matters: it marks a tuple with
one item, so the array has one axis containing four positions.

\hh{Create useful number sequences}

NumPy can create evenly spaced values without typing every number.

\code{c11_create_sequences}{python}{
    import numpy as np

    floor_numbers = np.arange(1, 6)
    sample_points_m = np.linspace(0.0, 10.0, num=5)
    blank_loads_kn = np.zeros(4)

    print(floor_numbers)
    print(sample_points_m)
    print(blank_loads_kn)
}{Create ranges, sample points, and zeros}

Use \c{arange()} when the step is the natural description. Use
\c{linspace()} when the number of samples, including both endpoints, is the
natural description.

\hh{Select values by position}

Array indices start at zero. Slices stop before the ending index, just as list
slices do.

\code{c11_index_and_slice}{python}{
    import numpy as np

    temperatures_c = np.array([19.5, 20.0, 21.5, 23.0, 22.0])

    print("First:", temperatures_c[0])
    print("Last:", temperatures_c[-1])
    print("Middle three:", temperatures_c[1:4])
}{Index and slice a one-dimensional array}

\begin{NoteBox}
\textbf{Stop and predict.} What values will \c{temperatures_c[::2]} select?
Check your answer by adding one print statement.
\end{NoteBox}

\hh{Calculate with a whole array}

Arithmetic operators work element by element. This is called a
\emph{vectorized calculation}. It expresses what should happen to the whole
collection without writing the loop yourself.

\code{c11_vectorized_arithmetic}{python}{
    import numpy as np

    widths_m = np.array([3.0, 4.0, 5.0, 3.5])
    lengths_m = np.array([6.0, 5.0, 4.0, 8.0])

    areas_m2 = widths_m * lengths_m
    adjusted_m2 = areas_m2 * 1.10

    print("Areas:", areas_m2)
    print("With 10 percent allowance:", adjusted_m2)
}{Multiply matching arrays element by element}

The arrays must have compatible shapes. Here each width is paired with the
length at the same position.

\hh{Compare a list operation with an array operation}

The multiplication operator has different meanings for lists and arrays.

\code{c11_list_array_difference}{python}{
    import numpy as np

    list_values = [2, 4, 6]
    array_values = np.array([2, 4, 6])

    print(list_values * 2)
    print(array_values * 2)
}{See why data type matters during multiplication}

The list is repeated. Each array value is multiplied. Inspect the type of a
value before assuming what an operator will do.

\hh{Filter with a boolean mask}

A comparison against an array produces one Boolean value per element. Use that
Boolean array as a mask to keep matching values.

\code{c11_boolean_mask}{python}{
    import numpy as np

    room_areas_m2 = np.array([24.0, 18.5, 31.0, 22.5, 35.0])
    large_mask = room_areas_m2 >= 25.0

    print("Mask:", large_mask)
    print("Matching areas:", room_areas_m2[large_mask])
}{Build and apply a boolean mask}

Parentheses keep combined conditions readable. Use \c{\&} for element-wise
and, \c{|} for element-wise or, and \c{\~} for element-wise not.

\code{c11_combined_mask}{python}{
    import numpy as np

    room_areas_m2 = np.array([24.0, 18.5, 31.0, 22.5, 35.0])
    middle_mask = (room_areas_m2 >= 20.0) & (room_areas_m2 <= 30.0)

    print(room_areas_m2[middle_mask])
}{Combine two array conditions}

Do not use Python's \c{and} or \c{or} to combine array comparisons. Those
operators expect one Boolean result, not one result per element.

\hh{Summarize numerical data}

Aggregation methods reduce many values to a smaller result.

\code{c11_basic_aggregates}{python}{
    import numpy as np

    room_areas_m2 = np.array([24.0, 18.5, 31.0, 22.5, 35.0])

    print("Count:", room_areas_m2.size)
    print("Total:", room_areas_m2.sum())
    print("Mean:", room_areas_m2.mean())
    print("Minimum:", room_areas_m2.min())
    print("Maximum:", room_areas_m2.max())
}{Calculate common summaries}

Choose the summary that answers the question. A mean is not a replacement for
looking at the minimum, maximum, and number of observations.

\hh{Represent a table with a two-dimensional array}

A nested list creates a two-dimensional array. The first axis identifies rows;
the second identifies columns.

\code{c11_two_dimensional_array}{python}{
    import numpy as np

    floor_areas_m2 = np.array([
        [24.0, 18.0, 30.0],
        [25.0, 19.0, 31.0],
        [23.0, 20.0, 32.0],
    ])

    print(floor_areas_m2)
    print("Shape:", floor_areas_m2.shape)
}{Create a three-row by three-column array}

The shape \c{(3, 3)} means three rows and three columns. NumPy arrays must be
rectangular: every row needs the same number of values.

\hh{Select rows, columns, and cells}

For a two-dimensional array, write the row selection first and the column
selection second.

\code{c11_select_2d}{python}{
    import numpy as np

    data = np.array([
        [24.0, 18.0, 30.0],
        [25.0, 19.0, 31.0],
        [23.0, 20.0, 32.0],
    ])

    print("One cell:", data[1, 2])
    print("First row:", data[0, :])
    print("Second column:", data[:, 1])
    print("Upper-left block:\n", data[:2, :2])
}{Select from a two-dimensional array}

The colon means all positions on that axis. Thus \c{data[:, 1]} selects every
row from column 1.

\hh{Summarize along an axis}

Without an \c{axis} argument, an aggregation uses every value. With an axis,
the aggregation removes that axis and keeps the other one.

\code{c11_axis_aggregates}{python}{
    import numpy as np

    floor_areas_m2 = np.array([
        [24.0, 18.0, 30.0],
        [25.0, 19.0, 31.0],
        [23.0, 20.0, 32.0],
    ])

    totals_by_column = floor_areas_m2.sum(axis=0)
    totals_by_row = floor_areas_m2.sum(axis=1)

    print("Column totals:", totals_by_column)
    print("Row totals:", totals_by_row)
}{Compare column and row summaries}

A useful check is the result's length: \c{axis=0} produces one result per
column, while \c{axis=1} produces one result per row.

\begin{NoteBox}
\textbf{Stop and predict.} For an array with shape \c{(4, 3)}, how many values
will \c{mean(axis=0)} return? How many will \c{mean(axis=1)} return?
\end{NoteBox}

\hh{Use simple broadcasting}

Broadcasting lets NumPy combine compatible shapes. A single number can apply
to every element. A one-dimensional array can also apply across matching
columns of a two-dimensional array.

\code{c11_broadcast_scalar}{python}{
    import numpy as np

    lengths_m = np.array([2.0, 3.5, 5.0])
    lengths_mm = lengths_m * 1000

    print(lengths_mm)
}{Broadcast one conversion factor across an array}

\code{c11_broadcast_columns}{python}{
    import numpy as np

    base_values = np.array([
        [10.0, 20.0, 30.0],
        [12.0, 22.0, 32.0],
    ])
    column_offsets = np.array([0.5, 1.0, 1.5])

    adjusted = base_values + column_offsets
    print(adjusted)
}{Broadcast three offsets across matching columns}

Compare shapes from right to left. Dimensions are compatible when they are
equal or one of them is 1. Begin with scalar broadcasting and matching columns;
more complex shape changes can wait until these cases feel natural.

\hh{Reshape without changing the values}

Reshaping changes the arrangement, not the total number of elements.

\code{c11_reshape_array}{python}{
    import numpy as np

    readings = np.arange(1, 13)
    table = readings.reshape(3, 4)

    print(table)
    print("Original size:", readings.size)
    print("New shape:", table.shape)
}{Reshape twelve values into three rows}

The requested dimensions must multiply to the original size. Twelve values can
form a \c{3} by \c{4} array, but not a \c{5} by \c{3} array.

\hh{Know when a slice shares data}

Basic NumPy slices usually return a \emph{view}. A view and the original array
share underlying data. This differs from a list slice, which creates a new
list.

\code{c11_slice_view}{python}{
    import numpy as np

    original = np.array([10.0, 20.0, 30.0, 40.0])
    selected = original[1:3]
    selected[0] = 99.0

    print("Selected:", selected)
    print("Original:", original)
}{Observe how a NumPy slice can change its source}

When the selected values must be independent, call \c{copy()} before changing
them.

\code{c11_safe_copy}{python}{
    import numpy as np

    original = np.array([10.0, 20.0, 30.0, 40.0])
    selected = original[1:3].copy()
    selected[0] = 99.0

    print("Selected copy:", selected)
    print("Original:", original)
}{Make selected array data independent}

Boolean indexing creates a copy, but relying on memorized exceptions is risky.
If independent data is part of your intention, make that intention visible with
\c{copy()}.

\hh{Common mistakes}

\begin{itemize}
    \item \textbf{Forgetting the import}: call NumPy functions through \c{np}
    after \c{import numpy as np}.
    \item \textbf{Expecting list behavior}: list multiplication repeats a
    list; array multiplication changes each value.
    \item \textbf{Mixing incompatible shapes}: inspect \c{shape} before
    combining arrays.
    \item \textbf{Using \c{and} or \c{or} with masks}: use parenthesized
    comparisons joined by \c{\&} or \c{|}.
    \item \textbf{Confusing axes}: connect the output length to the rows or
    columns that remain.
    \item \textbf{Changing a view accidentally}: call \c{copy()} before
    modifying selected data that must be independent.
\end{itemize}

\hh{Practice ladder}

\hhh{Level 0 -- Read and predict}

Without running the code, predict the two printed arrays.

\code{c11_practice_level0}{python}{
    import numpy as np

    values = np.array([2, 4, 6, 8])
    print(values + 1)
    print(values[1:3])
}{Predict vectorized addition and slicing}

\hhh{Level 1 -- Modify}

Change \c{threshold} to \c{25.0}. Then change the comparison so values equal to
the threshold are included.

\code{c11_practice_level1}{python}{
    import numpy as np

    areas_m2 = np.array([18.0, 24.0, 25.0, 31.0])
    threshold = 20.0
    print(areas_m2[areas_m2 > threshold])
}{Modify a filtering condition}

\hhh{Level 2 -- Complete}

Create a \c{2} by \c{3} array named \c{loads_kn}. Print its shape, the total
of each row, and the maximum value in the whole array.

\hhh{Level 3 -- Apply}

The array below contains four widths and four lengths. Calculate every area,
select areas of at least \c{20.0} square metres, and print their mean.

\begin{verbatim}
widths_m = np.array([3.0, 4.5, 5.0, 2.5])
lengths_m = np.array([6.0, 5.0, 4.5, 7.0])
\end{verbatim}

\textbf{Hints}
\begin{itemize}
    \item Multiply the two arrays element by element.
    \item Build a Boolean mask from the calculated areas.
    \item Apply the mask before calling \c{mean()}.
\end{itemize}

\hhh{Level 4 -- Challenge}

Write a function named \c{column_summary(data)}. It should return three arrays:
the minimum, mean, and maximum of each column. Test it with a three-row array.

\hh{Practice solutions}

\textbf{Level 0:} The outputs are \c{[3 5 7 9]} and \c{[4 6]}.

\textbf{Level 2}

\code{c11_solution_level2}{python}{
    import numpy as np

    loads_kn = np.array([
        [12.0, 15.0, 11.0],
        [14.0, 13.0, 16.0],
    ])

    print("Shape:", loads_kn.shape)
    print("Row totals:", loads_kn.sum(axis=1))
    print("Maximum:", loads_kn.max())
}{One solution to Level 2}

\textbf{Level 3}

\code{c11_solution_level3}{python}{
    import numpy as np

    widths_m = np.array([3.0, 4.5, 5.0, 2.5])
    lengths_m = np.array([6.0, 5.0, 4.5, 7.0])
    areas_m2 = widths_m * lengths_m
    selected_m2 = areas_m2[areas_m2 >= 20.0]

    print("Areas:", areas_m2)
    print("Selected:", selected_m2)
    print("Selected mean:", selected_m2.mean())
}{One solution to Level 3}

\textbf{Level 4}

\code{c11_solution_level4}{python}{
    import numpy as np

    def column_summary(data):
        minimum = data.min(axis=0)
        mean = data.mean(axis=0)
        maximum = data.max(axis=0)
        return minimum, mean, maximum

    readings = np.array([
        [18.0, 24.0, 30.0],
        [20.0, 22.0, 34.0],
        [19.0, 26.0, 32.0],
    ])

    minimum, mean, maximum = column_summary(readings)
    print("Minimum:", minimum)
    print("Mean:", mean)
    print("Maximum:", maximum)
}{One solution to Level 4}

\hh{Part IV checkpoint}

\begin{enumerate}
    \item Why are NumPy arrays useful for numerical collections?
    \item What does an array's \c{shape} describe?
    \item How does \c{size} differ from \c{ndim}?
    \item What does vectorized arithmetic do?
    \item What is a Boolean mask?
    \item In a two-dimensional array, which index is written first?
    \item What does \c{sum(axis=0)} return for a table?
    \item When can two shapes broadcast together?
    \item What must stay unchanged during \c{reshape()}?
    \item Why might a sliced array need \c{copy()}?
\end{enumerate}

\hh{Checkpoint answers}

\begin{enumerate}
    \item They store regular numerical data and support concise whole-array calculations.
    \item The number of positions along each axis.
    \item \c{size} counts all elements; \c{ndim} counts axes.
    \item It applies an operation across array elements without an explicit Python loop.
    \item A Boolean array used to select positions where a condition is true.
    \item The row index, followed by the column index.
    \item One total for each column.
    \item Corresponding dimensions, compared from the right, must be equal or one must be 1.
    \item The total number of elements.
    \item A basic slice may share data with the original array.
\end{enumerate}

\begin{tcolorbox}[title=Chapter 11 summary,colframe=orange,colback=white,arc=2mm]
NumPy arrays give numerical data a clear shape and support element-wise
calculations, filtering, and aggregation. Inspect shapes before combining
arrays, connect axis choices to the output you expect, and copy selected data
when it must be independent. The next chapter adds row and column labels with
pandas.
\end{tcolorbox}
